\documentclass[12pt,a4paper]{exam}
\usepackage[utf8]{inputenc}
\usepackage{amsmath,amssymb,amsthm}
\usepackage[margin=1in]{geometry}
\usepackage{graphicx}
\usepackage[colorlinks=true]{hyperref}
\pagestyle{headandfoot}
\firstpageheader{MIT OpenCourseWare}{}{\textbf{EXTREMELY HARD -- MIT LEVEL}}
\runningheader{MIT OpenCourseWare}{}{\textbf{EXTREMELY HARD -- MIT LEVEL}}
\firstpagefooter{}{Page \thepage}{}
\runningfooter{}{Page \thepage}{}

\begin{document}

\begin{center}
{\LARGE \textbf{MIT OpenCourseWare}}\\14.01SC\\Principles of Microeconomics
\vspace{0.5cm}
{14.01SC} --- {Principles of Microeconomics}
\vspace{0.3cm}
May 2026
\vspace{0.3cm}
\textbf{EXTREMELY HARD -- MIT LEVEL}
\end{center}

\begin{center}
\textbf{Time Limit: 3 hours}
\end{center}

\begin{center}
\textbf{Total Marks: 100}
\end{center}

\vspace{0.5cm}

\begin{questions}

\question[4]
Construct an original proof-level problem involving Supply and Demand for 14.01SC (Principles of Microeconomics) and provide a complete solution. Your problem should be at the level of a graduate qualifying exam.

\vspace{0.3cm}

\question[5]
Analyze the limitations of standard approaches to Consumer Theory when applied to edge cases in 14.01SC. Construct explicit counterexamples showing where intuitive reasoning fails.

\vspace{0.3cm}

\question[5]
Analyze the limitations of standard approaches to Producer Theory when applied to edge cases in 14.01SC. Construct explicit counterexamples showing where intuitive reasoning fails.

\vspace{0.3cm}

\question[4]
Analyze the limitations of standard approaches to Market Structure when applied to edge cases in 14.01SC. Construct explicit counterexamples showing where intuitive reasoning fails.

\vspace{0.3cm}

\question[3]
Prove the fundamental theorem concerning Game Theory in the context of 14.01SC (Principles of Microeconomics). State the theorem precisely, provide a complete proof, and discuss three important corollaries.

\vspace{0.3cm}

\question[4]
Analyze the limitations of standard approaches to Welfare Economics when applied to edge cases in 14.01SC. Construct explicit counterexamples showing where intuitive reasoning fails.

\vspace{0.3cm}

\question[4]
Analyze the limitations of standard approaches to Supply and Demand when applied to edge cases in 14.01SC. Construct explicit counterexamples showing where intuitive reasoning fails.

\vspace{0.3cm}

\question[7]
Construct an original proof-level problem involving Consumer Theory for 14.01SC (Principles of Microeconomics) and provide a complete solution. Your problem should be at the level of a graduate qualifying exam.

\vspace{0.3cm}

\question[4]
Construct an original proof-level problem involving Producer Theory for 14.01SC (Principles of Microeconomics) and provide a complete solution. Your problem should be at the level of a graduate qualifying exam.

\vspace{0.3cm}

\question[3]
Prove the fundamental theorem concerning Market Structure in the context of 14.01SC (Principles of Microeconomics). State the theorem precisely, provide a complete proof, and discuss three important corollaries.

\vspace{0.3cm}

\question[3]
Construct an original proof-level problem involving Game Theory for 14.01SC (Principles of Microeconomics) and provide a complete solution. Your problem should be at the level of a graduate qualifying exam.

\vspace{0.3cm}

\question[4]
Prove the fundamental theorem concerning Welfare Economics in the context of 14.01SC (Principles of Microeconomics). State the theorem precisely, provide a complete proof, and discuss three important corollaries.

\vspace{0.3cm}

\question[3]
Prove the fundamental theorem concerning Supply and Demand in the context of 14.01SC (Principles of Microeconomics). State the theorem precisely, provide a complete proof, and discuss three important corollaries.

\vspace{0.3cm}

\question[5]
Prove the fundamental theorem concerning Consumer Theory in the context of 14.01SC (Principles of Microeconomics). State the theorem precisely, provide a complete proof, and discuss three important corollaries.

\vspace{0.3cm}

\question[4]
Construct an original proof-level problem involving Producer Theory for 14.01SC (Principles of Microeconomics) and provide a complete solution. Your problem should be at the level of a graduate qualifying exam.

\vspace{0.3cm}

\question[4]
Prove the fundamental theorem concerning Market Structure in the context of 14.01SC (Principles of Microeconomics). State the theorem precisely, provide a complete proof, and discuss three important corollaries.

\vspace{0.3cm}

\question[3]
Prove the fundamental theorem concerning Game Theory in the context of 14.01SC (Principles of Microeconomics). State the theorem precisely, provide a complete proof, and discuss three important corollaries.

\vspace{0.3cm}

\question[3]
Construct an original proof-level problem involving Welfare Economics for 14.01SC (Principles of Microeconomics) and provide a complete solution. Your problem should be at the level of a graduate qualifying exam.

\vspace{0.3cm}

\question[3]
Construct an original proof-level problem involving Supply and Demand for 14.01SC (Principles of Microeconomics) and provide a complete solution. Your problem should be at the level of a graduate qualifying exam.

\vspace{0.3cm}

\question[3]
Analyze the limitations of standard approaches to Consumer Theory when applied to edge cases in 14.01SC. Construct explicit counterexamples showing where intuitive reasoning fails.

\vspace{0.3cm}

\question[6]
Construct an original proof-level problem involving Producer Theory for 14.01SC (Principles of Microeconomics) and provide a complete solution. Your problem should be at the level of a graduate qualifying exam.

\vspace{0.3cm}

\question[4]
Prove the fundamental theorem concerning Market Structure in the context of 14.01SC (Principles of Microeconomics). State the theorem precisely, provide a complete proof, and discuss three important corollaries.

\vspace{0.3cm}

\question[2]
Analyze the limitations of standard approaches to Game Theory when applied to edge cases in 14.01SC. Construct explicit counterexamples showing where intuitive reasoning fails.

\vspace{0.3cm}

\question[5]
Construct an original proof-level problem involving Welfare Economics for 14.01SC (Principles of Microeconomics) and provide a complete solution. Your problem should be at the level of a graduate qualifying exam.

\vspace{0.3cm}

\question[5]
Analyze the limitations of standard approaches to Supply and Demand when applied to edge cases in 14.01SC. Construct explicit counterexamples showing where intuitive reasoning fails.

\vspace{0.3cm}

\end{questions}

\newpage
\section*{Solutions}

\textbf{Solution 1:}

Solution approach: First recall the definitions and key theorems related to Supply and Demand from 14.01SC. Then apply the appropriate techniques to construct a rigorous argument. The solution must be self-contained and reference only the material from the course.

\vspace{0.5cm}

\textbf{Solution 2:}

The edge case analysis reveals the subtle assumptions underlying standard treatments of Consumer Theory. By constructing explicit counterexamples, we see that the usual hypotheses cannot be weakened without loss of the theorem's conclusion.

\vspace{0.5cm}

\textbf{Solution 3:}

The edge case analysis reveals the subtle assumptions underlying standard treatments of Producer Theory. By constructing explicit counterexamples, we see that the usual hypotheses cannot be weakened without loss of the theorem's conclusion.

\vspace{0.5cm}

\textbf{Solution 4:}

The edge case analysis reveals the subtle assumptions underlying standard treatments of Market Structure. By constructing explicit counterexamples, we see that the usual hypotheses cannot be weakened without loss of the theorem's conclusion.

\vspace{0.5cm}

\textbf{Solution 5:}

The proof follows from the definitions and properties established in 14.01SC. The key insight is to apply the relevant theorem and verify the hypotheses hold. Detailed solution depends on the specific formulation of the theorem.

\vspace{0.5cm}

\textbf{Solution 6:}

The edge case analysis reveals the subtle assumptions underlying standard treatments of Welfare Economics. By constructing explicit counterexamples, we see that the usual hypotheses cannot be weakened without loss of the theorem's conclusion.

\vspace{0.5cm}

\textbf{Solution 7:}

The edge case analysis reveals the subtle assumptions underlying standard treatments of Supply and Demand. By constructing explicit counterexamples, we see that the usual hypotheses cannot be weakened without loss of the theorem's conclusion.

\vspace{0.5cm}

\textbf{Solution 8:}

Solution approach: First recall the definitions and key theorems related to Consumer Theory from 14.01SC. Then apply the appropriate techniques to construct a rigorous argument. The solution must be self-contained and reference only the material from the course.

\vspace{0.5cm}

\textbf{Solution 9:}

Solution approach: First recall the definitions and key theorems related to Producer Theory from 14.01SC. Then apply the appropriate techniques to construct a rigorous argument. The solution must be self-contained and reference only the material from the course.

\vspace{0.5cm}

\textbf{Solution 10:}

The proof follows from the definitions and properties established in 14.01SC. The key insight is to apply the relevant theorem and verify the hypotheses hold. Detailed solution depends on the specific formulation of the theorem.

\vspace{0.5cm}

\textbf{Solution 11:}

Solution approach: First recall the definitions and key theorems related to Game Theory from 14.01SC. Then apply the appropriate techniques to construct a rigorous argument. The solution must be self-contained and reference only the material from the course.

\vspace{0.5cm}

\textbf{Solution 12:}

The proof follows from the definitions and properties established in 14.01SC. The key insight is to apply the relevant theorem and verify the hypotheses hold. Detailed solution depends on the specific formulation of the theorem.

\vspace{0.5cm}

\textbf{Solution 13:}

The proof follows from the definitions and properties established in 14.01SC. The key insight is to apply the relevant theorem and verify the hypotheses hold. Detailed solution depends on the specific formulation of the theorem.

\vspace{0.5cm}

\textbf{Solution 14:}

The proof follows from the definitions and properties established in 14.01SC. The key insight is to apply the relevant theorem and verify the hypotheses hold. Detailed solution depends on the specific formulation of the theorem.

\vspace{0.5cm}

\textbf{Solution 15:}

Solution approach: First recall the definitions and key theorems related to Producer Theory from 14.01SC. Then apply the appropriate techniques to construct a rigorous argument. The solution must be self-contained and reference only the material from the course.

\vspace{0.5cm}

\textbf{Solution 16:}

The proof follows from the definitions and properties established in 14.01SC. The key insight is to apply the relevant theorem and verify the hypotheses hold. Detailed solution depends on the specific formulation of the theorem.

\vspace{0.5cm}

\textbf{Solution 17:}

The proof follows from the definitions and properties established in 14.01SC. The key insight is to apply the relevant theorem and verify the hypotheses hold. Detailed solution depends on the specific formulation of the theorem.

\vspace{0.5cm}

\textbf{Solution 18:}

Solution approach: First recall the definitions and key theorems related to Welfare Economics from 14.01SC. Then apply the appropriate techniques to construct a rigorous argument. The solution must be self-contained and reference only the material from the course.

\vspace{0.5cm}

\textbf{Solution 19:}

Solution approach: First recall the definitions and key theorems related to Supply and Demand from 14.01SC. Then apply the appropriate techniques to construct a rigorous argument. The solution must be self-contained and reference only the material from the course.

\vspace{0.5cm}

\textbf{Solution 20:}

The edge case analysis reveals the subtle assumptions underlying standard treatments of Consumer Theory. By constructing explicit counterexamples, we see that the usual hypotheses cannot be weakened without loss of the theorem's conclusion.

\vspace{0.5cm}

\textbf{Solution 21:}

Solution approach: First recall the definitions and key theorems related to Producer Theory from 14.01SC. Then apply the appropriate techniques to construct a rigorous argument. The solution must be self-contained and reference only the material from the course.

\vspace{0.5cm}

\textbf{Solution 22:}

The proof follows from the definitions and properties established in 14.01SC. The key insight is to apply the relevant theorem and verify the hypotheses hold. Detailed solution depends on the specific formulation of the theorem.

\vspace{0.5cm}

\textbf{Solution 23:}

The edge case analysis reveals the subtle assumptions underlying standard treatments of Game Theory. By constructing explicit counterexamples, we see that the usual hypotheses cannot be weakened without loss of the theorem's conclusion.

\vspace{0.5cm}

\textbf{Solution 24:}

Solution approach: First recall the definitions and key theorems related to Welfare Economics from 14.01SC. Then apply the appropriate techniques to construct a rigorous argument. The solution must be self-contained and reference only the material from the course.

\vspace{0.5cm}

\textbf{Solution 25:}

The edge case analysis reveals the subtle assumptions underlying standard treatments of Supply and Demand. By constructing explicit counterexamples, we see that the usual hypotheses cannot be weakened without loss of the theorem's conclusion.

\vspace{0.5cm}

\end{document}